\documentclass[11pt]{article}
\usepackage[margin=1.2in, top=1in, bottom=1in]{geometry}
\usepackage{amsmath, amssymb, amsthm}
\usepackage{microtype}
\usepackage{parskip}
\usepackage[T1]{fontenc}
\usepackage{lmodern}

\setlength{\parindent}{0pt}
\setlength{\parskip}{6pt}

\begin{document}

\begin{center}
  {\large\bfseries RetroDiff: Training Objective Derivation}
\end{center}

\medskip

\noindent
Let $x$ denote a clean data sample and $x_t$ its noised counterpart at diffusion step $t \in \{1,\ldots,T\}$.
The forward process $q(x_{1:T}|x)$ gradually corrupts $x$; the model $p_\theta$ learns to reverse it.
Applying Jensen's inequality to $\log p_\theta(x)$ yields the variational lower bound (ELBO):

\begin{align}
  \log p_\theta(x)
    &\;\geq\;
    \mathbb{E}_{q(x_{1:T}|x)}\!\left[\log\frac{p_\theta(x_{0:T})}{q(x_{1:T}|x)}\right]
    \;=\; -\mathcal{L}_{\mathrm{ELBO}}
  \label{eq:elbo}
\end{align}

\noindent
The bound in \eqref{eq:elbo} is tight when the variational posterior matches the true reverse process.
Expanding the joint $p_\theta(x_{0:T})$ and the Markov structure of $q$ decomposes the ELBO into a
reconstruction term at $t{=}1$ and a sum of per-step KL divergences:

\begin{align}
  \mathcal{L}_{\mathrm{ELBO}}
    &\;=\;
    \underbrace{\mathbb{E}_q\!\left[-\log p_\theta(x_0|x_1)\right]}_{\text{reconstruction}}
    \;+\;
    \sum_{t=2}^{T}
    \underbrace{\mathrm{KL}\!\left(q(x_{t-1}|x_t,x)\;\|\;p_\theta(x_{t-1}|x_t)\right)}_{\text{step-}t\text{ denoising}}
  \label{eq:decomp}
\end{align}

\noindent
Each KL term in \eqref{eq:decomp} is tractable because the forward posterior
$q(x_{t-1}|x_t,x)$ is Gaussian with a closed-form mean and variance.
\textbf{RetroDiff}'s key contribution is to condition the entire objective on a set
$R(x)$ of $k$ retrieved exemplars drawn from the training corpus at inference time,
giving the retrieval-augmented training objective:

\begin{align}
  \mathcal{L}_{\mathrm{RetroDiff}}(\theta)
    &\;=\;
    \mathbb{E}_{x \sim p_{\mathrm{data}}}\!\left[
      \mathcal{L}_{\mathrm{ELBO}}\!\left(x \;\big|\; R(x)\right)
    \right]
  \label{eq:retro}
\end{align}

\noindent
Minimising \eqref{eq:retro} trains the reverse process to leverage retrieved neighbours,
reducing the effective denoising burden and improving sample quality without altering
the diffusion schedule.

\end{document}
